\subsection{D-ACT Denial Attack Simulator}
\label{app:dact_denial_attack_simulator}

The supplementary web artefact includes a Denial Attack Simulator page. The simulator is a purely visual educational tool designed to explain the relationship between the D-ACT taxonomy, denial attack classes, active conditions, victim resource domains, and mitigation strategies. It does not generate packets, replay packet captures, create HTTP requests, or interact with external infrastructure. Instead, it animates abstract packet and resource-flow indicators between an attacker-side visual and a victim-side visual.

The simulator supports the nine primary taxonomy classes used in the D-ACT evaluation corpus: DoS, DDoS, LDoS, LDDoS, \(\text{EDoS}_{S}\), \(\text{EDoS}_{D}\), DoW, DDoW, and DoA. For each class, the simulator highlights the active condition set from \(C_0\) to \(C_6\). For example, DoS activates \(C_0\) and \(C_1\), DDoS activates \(C_0\) and \(C_2\), LDoS activates \(C_0\), \(C_1\), and \(C_3\), and DoA activates \(C_0\) and \(C_6\). Cloud-targeted, serverless-targeted, and AI-targeted classes add visual resource stages corresponding to \(C_4\), \(C_5\), and \(C_6\), respectively.

The simulator also integrates the mitigation catalogue used by the D-ACT web interface. When an attack class is selected, the side panel displays the relevant mitigation strategy, including immediate triage, technical controls, and monitoring or follow-up actions. Technical controls are interactive. Selecting a control changes the visual state of the simulation to represent the defensive effect, such as blocking, throttling, traffic filtering, endpoint hardening, cloud cost guardrails, serverless invocation caps, AI token limits, or monitoring overlays. The simulator therefore functions as a non-operational explanatory artefact that connects taxonomy theory, classifier output, and defensive response in a way that is accessible to non-technical users.
